\documentclass[a4paper]{article}

\def\npart{II}

\def\ntitle{Algebraic Topology}
\def\nlecturer{H.\ Wilton}

\def\nterm{Michaelmas}
\def\nyear{2018}

\ifx \nauthor\undefined
  \def\nauthor{Qiangru Kuang}
\else
\fi

\ifx \ntitle\undefined
  \def\ntitle{Template}
\else
\fi

\ifx \nauthoremail\undefined
  \def\nauthoremail{qk206@cam.ac.uk}
\else
\fi

\ifx \ndate\undefined
  \def\ndate{\today}
\else
\fi

\title{\ntitle}
\author{\nauthor}
\date{\ndate}

%\usepackage{microtype}
\usepackage{mathtools}
\usepackage{amsthm}
\usepackage{stmaryrd}%symbols used so far: \mapsfrom
\usepackage{empheq}
\usepackage{amssymb}
\let\mathbbalt\mathbb
\let\pitchforkold\pitchfork
\usepackage{unicode-math}
\let\mathbb\mathbbalt%reset to original \mathbb
\let\pitchfork\pitchforkold

\usepackage{imakeidx}
\makeindex[intoc]

%to address the problem that Latin modern doesn't have unicode support for setminus
%https://tex.stackexchange.com/a/55205/26707
\AtBeginDocument{\renewcommand*{\setminus}{\mathbin{\backslash}}}
\AtBeginDocument{\renewcommand*{\models}{\vDash}}%for \vDash is same size as \vdash but orginal \models is larger
\AtBeginDocument{\let\Re\relax}
\AtBeginDocument{\let\Im\relax}
\AtBeginDocument{\DeclareMathOperator{\Re}{Re}}
\AtBeginDocument{\DeclareMathOperator{\Im}{Im}}
\AtBeginDocument{\let\div\relax}
\AtBeginDocument{\DeclareMathOperator{\div}{div}}

\usepackage{tikz}
\usetikzlibrary{automata,positioning}
\usepackage{pgfplots}
%some preset styles
\pgfplotsset{compat=1.15}
\pgfplotsset{centre/.append style={axis x line=middle, axis y line=middle, xlabel={$x$}, ylabel={$y$}, axis equal}}
\usepackage{tikz-cd}
\usepackage{graphicx}
\usepackage{newunicodechar}

\usepackage{fancyhdr}

\fancypagestyle{mypagestyle}{
    \fancyhf{}
    \lhead{\emph{\nouppercase{\leftmark}}}
    \rhead{}
    \cfoot{\thepage}
}
\pagestyle{mypagestyle}

\usepackage{titlesec}
\newcommand{\sectionbreak}{\clearpage} % clear page after each section
\usepackage[perpage]{footmisc}
\usepackage{blindtext}

%\reallywidehat
%https://tex.stackexchange.com/a/101136/26707
\usepackage{scalerel,stackengine}
\stackMath
\newcommand\reallywidehat[1]{%
\savestack{\tmpbox}{\stretchto{%
  \scaleto{%
    \scalerel*[\widthof{\ensuremath{#1}}]{\kern-.6pt\bigwedge\kern-.6pt}%
    {\rule[-\textheight/2]{1ex}{\textheight}}%WIDTH-LIMITED BIG WEDGE
  }{\textheight}% 
}{0.5ex}}%
\stackon[1pt]{#1}{\tmpbox}%
}

%\usepackage{braket}
\usepackage{thmtools}%restate theorem
\usepackage{hyperref}

% https://en.wikibooks.org/wiki/LaTeX/Hyperlinks
\hypersetup{
    %bookmarks=true,
    unicode=true,
    pdftitle={\ntitle},
    pdfauthor={\nauthor},
    pdfsubject={Mathematics},
    pdfcreator={\nauthor},
    pdfproducer={\nauthor},
    pdfkeywords={math maths \ntitle},
    colorlinks=true,
    linkcolor={red!50!black},
    citecolor={blue!50!black},
    urlcolor={blue!80!black}
}

\usepackage{cleveref}



% TODO: mdframed often gives bad breaks that cause empty lines. Would like to switch to tcolorbox.
% The current workaround is to set innerbottommargin=0pt.

%\usepackage[theorems]{tcolorbox}





\usepackage[framemethod=tikz]{mdframed}
\mdfdefinestyle{leftbar}{
  %nobreak=true, %dirty hack
  linewidth=1.5pt,
  linecolor=gray,
  hidealllines=true,
  leftline=true,
  leftmargin=0pt,
  innerleftmargin=5pt,
  innerrightmargin=10pt,
  innertopmargin=-5pt,
  % innerbottommargin=5pt, % original
  innerbottommargin=0pt, % temporary hack 
}
%\newmdtheoremenv[style=leftbar]{theorem}{Theorem}[section]
%\newmdtheoremenv[style=leftbar]{proposition}[theorem]{proposition}
%\newmdtheoremenv[style=leftbar]{lemma}[theorem]{Lemma}
%\newmdtheoremenv[style=leftbar]{corollary}[theorem]{corollary}

\newtheorem{theorem}{Theorem}[section]
\newtheorem{proposition}[theorem]{Proposition}
\newtheorem{lemma}[theorem]{Lemma}
\newtheorem{corollary}[theorem]{Corollary}
\newtheorem{axiom}[theorem]{Axiom}
\newtheorem*{axiom*}{Axiom}

\surroundwithmdframed[style=leftbar]{theorem}
\surroundwithmdframed[style=leftbar]{proposition}
\surroundwithmdframed[style=leftbar]{lemma}
\surroundwithmdframed[style=leftbar]{corollary}
\surroundwithmdframed[style=leftbar]{axiom}
\surroundwithmdframed[style=leftbar]{axiom*}

\theoremstyle{definition}

\newtheorem*{definition}{Definition}
\surroundwithmdframed[style=leftbar]{definition}

\newtheorem*{slogan}{Slogan}
\newtheorem*{eg}{Example}
\newtheorem*{ex}{Exercise}
\newtheorem*{remark}{Remark}
\newtheorem*{notation}{Notation}
\newtheorem*{convention}{Convention}
\newtheorem*{assumption}{Assumption}
\newtheorem*{question}{Question}
\newtheorem*{answer}{Answer}
\newtheorem*{note}{Note}
\newtheorem*{application}{Application}

%operator macros

%basic
\DeclareMathOperator{\lcm}{lcm}

%matrix
\DeclareMathOperator{\tr}{tr}
\DeclareMathOperator{\Tr}{Tr}
\DeclareMathOperator{\adj}{adj}

%algebra
\DeclareMathOperator{\Hom}{Hom}
\DeclareMathOperator{\End}{End}
\DeclareMathOperator{\id}{id}
\DeclareMathOperator{\im}{im}
\DeclarePairedDelimiter{\generation}{\langle}{\rangle}

%groups
\DeclareMathOperator{\sym}{Sym}
\DeclareMathOperator{\sgn}{sgn}
\DeclareMathOperator{\inn}{Inn}
\DeclareMathOperator{\aut}{Aut}
\DeclareMathOperator{\GL}{GL}
\DeclareMathOperator{\SL}{SL}
\DeclareMathOperator{\PGL}{PGL}
\DeclareMathOperator{\PSL}{PSL}
\DeclareMathOperator{\SU}{SU}
\DeclareMathOperator{\UU}{U}
\DeclareMathOperator{\SO}{SO}
\DeclareMathOperator{\OO}{O}
\DeclareMathOperator{\PSU}{PSU}

%hyperbolic
\DeclareMathOperator{\sech}{sech}

%field, galois heory
\DeclareMathOperator{\ch}{ch}
\DeclareMathOperator{\gal}{Gal}
\DeclareMathOperator{\emb}{Emb}



%ceiling and floor
%https://tex.stackexchange.com/a/118217/26707
\DeclarePairedDelimiter\ceil{\lceil}{\rceil}
\DeclarePairedDelimiter\floor{\lfloor}{\rfloor}


\DeclarePairedDelimiter{\innerproduct}{\langle}{\rangle}

%\DeclarePairedDelimiterX{\norm}[1]{\lVert}{\rVert}{#1}
\DeclarePairedDelimiter{\norm}{\lVert}{\rVert}



%Dirac notation
%TODO: rewrite for variable number of arguments
\DeclarePairedDelimiterX{\braket}[2]{\langle}{\rangle}{#1 \delimsize\vert #2}
\DeclarePairedDelimiterX{\braketthree}[3]{\langle}{\rangle}{#1 \delimsize\vert #2 \delimsize\vert #3}

\DeclarePairedDelimiter{\bra}{\langle}{\rvert}
\DeclarePairedDelimiter{\ket}{\lvert}{\rangle}




%macros

%general

%divide, not divide
\newcommand*{\divides}{\mid}
\newcommand*{\ndivides}{\nmid}
%vector, i.e. mathbf
%https://tex.stackexchange.com/a/45746/26707
\newcommand*{\V}[1]{{\ensuremath{\symbf{#1}}}}
%closure
\newcommand*{\cl}[1]{\overline{#1}}
%conjugate
\newcommand*{\conj}[1]{\overline{#1}}
%set complement
\newcommand*{\stcomp}[1]{\overline{#1}}
\newcommand*{\compose}{\circ}
\newcommand*{\nto}{\nrightarrow}
\newcommand*{\p}{\partial}
%embed
\newcommand*{\embed}{\hookrightarrow}
%surjection
\newcommand*{\surj}{\twoheadrightarrow}
%power set
\newcommand*{\powerset}{\mathcal{P}}

%matrix
\newcommand*{\matrixring}{\mathcal{M}}

%groups
\newcommand*{\normal}{\trianglelefteq}
%rings
\newcommand*{\ideal}{\trianglelefteq}

%fields
\renewcommand*{\C}{{\mathbb{C}}}
\newcommand*{\R}{{\mathbb{R}}}
\newcommand*{\Q}{{\mathbb{Q}}}
\newcommand*{\Z}{{\mathbb{Z}}}
\newcommand*{\N}{{\mathbb{N}}}
\newcommand*{\F}{{\mathbb{F}}}
%not really but I think this belongs here
\newcommand*{\A}{{\mathbb{A}}}

%asymptotic
\newcommand*{\bigO}{O}
\newcommand*{\smallo}{o}

%probability
\newcommand*{\prob}{\mathbb{P}}
\newcommand*{\E}{\mathbb{E}}

%vector calculus
\newcommand*{\gradient}{\V \nabla}
\newcommand*{\divergence}{\gradient \cdot}
\newcommand*{\curl}{\gradient \cdot}

%logic
\newcommand*{\yields}{\vdash}
\newcommand*{\nyields}{\nvdash}

%differential geometry
\renewcommand*{\H}{\mathbb{H}}
\newcommand*{\transversal}{\pitchfork}
\renewcommand{\d}{\mathrm{d}} % exterior derivative

%number theory
\newcommand*{\legendre}[2]{\genfrac{(}{)}{}{}{#1}{#2}}%Legendre symbol


\DeclareMathOperator{\rel}{rel}

\begin{document}

\begin{titlepage}
  \begin{center}
    \includegraphics[width=0.6\textwidth]{logo.jpg}\par
    \vspace{1cm}
    {\scshape\huge Mathematics Tripos \par}
    \vspace{2cm}
    {\huge Part \npart \par}
    \vspace{0.6cm}
    {\Huge \bfseries \ntitle \par}
    \vspace{1.2cm}
    {\Large\nterm, \nyear \par}
    \vspace{2cm}
    
    {\large \emph{Lectures by } \par}
    \vspace{0.2cm}
    {\Large \scshape \nlecturer}
    
    \vspace{0.5cm}
    {\large \emph{Notes by }\par}
    \vspace{0.2cm}
    {\Large \scshape \href{mailto:\nauthoremail}{\nauthor}}
 \end{center}
\end{titlepage}

\tableofcontents

\setcounter{section}{-1}

\section{Introduction}

\begin{question}
  Is the Hopf link really linked? More formally, is there a homeomorphism \(\R^3 \to \R^3\) taking \(H\) to \(U\)?
\end{question}

\(H\) can be realised as \(S^1 \amalg S^1 \to \R^3\). For \(U\), we can consider \(S^1 \amalg S^1\) as boundary of \(D^1 \amalg D^1\) and the map extends to a map to discs.

So it makes sense to phrase the question as

\begin{question}
  Does the Hopf link \(\eta: S^1 \amalg S^1 \to \R^3\)extend to a map of discs?
\end{question}

This is an example of an \emph{extension problem}.

Here is another example. Define the \(n\)-sphere \(S^{n - 1} := \{x \in \R^n: \sum_{i = 1}^n x_i^2 = 1\}\), which sits inside \(D^n = \{x \in \R^n: \sum_{i = 1}^n x_i^2 = 1\}\). We can ask:

\begin{question}
  Does the identity map \(\id_{S^{n - 1}}: S^{n - 1} \to S^{n - 1}\) factor through \(D^n\)?
\end{question}

To gain some intuition, let's consider small \(n\). For \(n = 1\), \(S^0 = \{-1, 1\}\). The answer is no by Intermediate Value Theorem, or connectedness from topology. For \(n = 2\), this answer is again no by winding number argument. What about \(n \geq 3\)?

These problems are hard because we have to consider continuous maps between two spaces, which are in general very big and hard to compute. On the other hand, a comparable algebraic problem is

\begin{question}
  Does the map \(\id: \Z \to \Z\) factor through \(0\)?
\end{question}

Well that's much much easier!

\section{The fundamental group}

Throughout this course, ``maps'' mean continuous maps.

\subsection{Deforming maps and spaces}

\begin{definition}[homotopy]\index{homotopy}
  Let \(f_0, f_1: X \to Y\) be maps. A \emph{homotopy} between \(f_0\) and \(f_1\) is a map \(F: X \times [0, 1] \to X\) such that \(F(x, 0) = f_0(x)\) and \(F(x, 1) = f_1(x)\) for all \(x \in X\).

  If \(F\) exists, we say that \(f_0\) is \emph{homotopic} to \(f_1\) and write \(f_0 \simeq f_1\), or to emphasise the homotopy, \(f_0 \simeq_F f_1\).
\end{definition}

\begin{notation}
  \(I = [0, 1]\). We often write \(f_t(x) = F(x, t)\).
\end{notation}

\begin{eg}
  If \(Y\) is a convex region in \(\R^n\) then for any \(f_0, f_1: X \to Y\), the \emph{straightline homotopy} \(F(x, y) = t f_1(x) + (1 - t) f_0(x)\) is a homotopy \(f_0 \simeq f_1\).
\end{eg}

\begin{definition}[relative homotopy]
  If \(Z \subseteq X\) and \(F(z, t) = f_0(z) = f_1(z)\) for all \(z \in Z, t \in I\), then \(F\) is a \emph{homotopy relative to \(Z\)}, write \(f_0 \simeq_F f_1 \rel Z\).
\end{definition}

\begin{lemma}
  The relation \(\simeq\) (\(\rel Z\)) is an equivalence relation on maps \(X \to Y\).
\end{lemma}

\begin{proof}
  Reflexivity and symmetry are easy. For transitivity, suppose \(f_0 \simeq_{F_0} f_1 \simeq_{F_1} f_2\). Let
  \[
    F(x, t) =
    \begin{cases}
      F_0(x, 2t) & t \leq \frac{1}{2} \\
      F_1(x, 2t - 1) & t \geq \frac{1}{2}
    \end{cases}
  \]
  which is the homotopy we need.
\end{proof}

\begin{definition}[homotopy equivalence]\index{homotopy equivalence}
  \(f: X \to Y\) and \(g: Y \to X\) is a \emph{homotopy equivalence} if \(g \compose f \simeq \id_X\) and \(f \compose g \simeq \id_Y\). In this case we say \(X\) is homotopy equivalent to \(Y\) and write \(X \simeq Y\).
\end{definition}

\begin{eg}
  Let \(X = *\), the space with one point and \(Y = \R^n\). Let \(f: * \mapsto 0\), \(g\) be the unique map \(Y \to X\). Then \(g \compose f = \id_X\), and \(f \compose g = 0 \simeq \id_Y\) via the straightline homotopy. Therefore \(\R^n\) is homotopy equivalent to \(*\).
\end{eg}

\begin{definition}[contractible]\index{contractible}
  A space \(X\) is \emph{contractible} if \(X \simeq *\).
\end{definition}

\begin{eg}
  Let \(X = S^1, Y = \R^2 - \{0\}\). Let \(f: X \to Y\) be the natural inclusion nad \(g: Y \to X, x \mapsto \frac{x}{\norm x}\). Then
  \begin{align*}
    g \compose f &= \id_X \\
    f \compose g(x) &= \frac{x}{\norm x} \in \R^2
  \end{align*}
  Although \(Y\) is not convex, for all \(x, t\), straightline homotopy \(F(x, t)\) between \(f \compose g\) and \(\id_Y\) satisfies \(F(x, t) \neq 0\) so \(f \compose g \simeq_F \id_Y\). Thus \(X \simeq Y\).
\end{eg}

\begin{definition}[retract, deformation retract]\index{retract}\index{deformation retract}
  Let \(f: X \to Y\) and \(g: Y \to X\). If \(g \compose f = \id_X\) then \(X\) is a \emph{retract} of \(Y\).

  If in addition \(f \compose g \simeq \id_Y \rel f(X)\) then we say \(X\) is a \emph{deformation retract} of \(Y\).
\end{definition}

Note that whenever we have \(g \compose f = \id_X\), \(f\) is injective so we can think \(X\) as being embedded in \(Y\). Informally, \(Y\) is ``as complicated'' as \(X\).

\begin{lemma}
  Homotopy equivalence is an equivalence on topological spaces.
\end{lemma}

\begin{proof}
  Symmetry and reflexivity are obvious. For transitivity, consider
  \[
    \begin{tikzcd}
      X \ar[r, "f", shift left] & Y \ar[l, "g", shift left] \ar[r, "f", shift left] & Z \ar[l, "g", shift left]
    \end{tikzcd}
  \]
  Need to show that \(g \compose (g' \compose f') \compose f \simeq \id_X\) (and the other direction will follow similarly). By hypothesis \(g' \compose f' \simeq_{F'} \id_Y\). Now
  \[
    g(F'(f(x), t))
  \]
  is a homotopy
  \[
    g \compose g' \compose f' \compose f \simeq g \compose \id_Y \compose f = g \compose f \simeq \id_X.
  \]
\end{proof}

\subsection{The fundamental group}

\begin{definition}[path, loop]\index{path}\index{loop}
  A \emph{path} (from \(x_0\) to \(x_1\)) is a continuous map \(\gamma: I \to X\) (with \(\gamma(0) = x_0, \gamma(1) = x_1\)).

  A \emph{loop} (based at \(x_0\)) is a path from \(x_0\) to \(x_0\).
\end{definition}

\begin{definition}[homotopy of path]\index{homotopy of path}
  Let \(\gamma_0, \gamma_1\) be paths from \(x_0\) to \(x_1\). A \emph{homotopy (of path)} from \(\gamma_0\) to \(\gamma_1\) is a homotopy
  \[
    \gamma_0 \simeq_F \gamma_1 \rel \{0, 1\}.
  \]
\end{definition}

\begin{definition}[concatenation of path, constant path, inverse path]
  Let \(\gamma\) be a path from \(x\) to \(y\) and \(\delta\) a path from \(y\) to \(z\).
  \begin{enumerate}
  \item The \emph{concatenation} of \(\gamma\) and \(\delta\) is
    \[
      (\gamma \cdot \delta) (t) =
      \begin{cases}
        \gamma(2t) & t \leq \frac{1}{2} \\
        \delta(2t - 1) & t \geq \frac{1}{2}
      \end{cases}
    \]
  \item The \emph{constant} path (at \(x\)) is \(c_x(t) = x\).
  \item The \emph{inverse path} to \(\gamma\) is \(\overline \gamma(t) = \gamma(1 - t)\).
  \end{enumerate}
\end{definition}

\begin{theorem}[fundamental group]\index{fundamental group}
  Let \(x_0 \in X\). Let
  \[
    \pi_1(X, x_0) = \{\text{loops based at } x_0\} / \simeq.
  \]
  This has a group structure with
  \begin{itemize}
  \item \([\gamma][\delta] = [\gamma \cdot \delta]\),
  \item identity \([c_{x_0}]\),
  \item \([\gamma]^{-1} = [\overline \gamma]\).
  \end{itemize}

  We call \(\pi_1(X, x_0)\) the \emph{fundamental group} of \(X\) (based at \(x_0\)).
\end{theorem}

\begin{proof}
  To prove the theorem, we need to check that multiplication and inverses are well-defined and the group axioms are satisfied.

  \begin{lemma}
    If \(\gamma_0, \gamma_1\) are paths to \(y\) and \(\delta_0, \delta_1\) are paths from \(y\) and \(\gamma_0 \simeq \gamma_1, \delta_0 \simeq \delta_1\), then
    \[
      \gamma_0 \cdot \delta_0 \simeq \gamma_1 \cdot \delta_1.
    \]

    Also \(\overline \gamma_0 \simeq \overline \gamma_1\).
  \end{lemma}

  \begin{proof}
    We only show for concatenation. Inverses are similar. Let \(\gamma_0 \simeq_F \gamma_1, \delta_0 \simeq_G \delta_1\). (proof by picture) Algebraically, the homotopy is given by
    \[
      H(s, t) =
      \begin{cases}
        F(s, 2t) & t \leq \frac{1}{2} \\
        G(s, 2t - 1) & t \geq \frac{1}{2}
      \end{cases}
    \]
  \end{proof}

  Now we check that the group axioms are satisfied.

  \begin{lemma}\leavevmode
    \begin{enumerate}
    \item \((\alpha \cdot \beta) \cdot \gamma \simeq \alpha \cdot (\beta \cdot \gamma)\).
    \item \(\alpha \cdot c_x \simeq \alpha \simeq c_w \cdot \alpha\).
    \item \(\alpha \cdot \overline \alpha \simeq c_w\).
    \end{enumerate}
  \end{lemma}

  \begin{proof}
    We show \(1\). The other two are similar. Let
    \[
      \delta =
      \begin{cases}
        \alpha(3t) & t \leq \frac{1}{3} \\
        \beta(3t - 1) & \frac{1}{3} \leq t \leq \frac{2}{3} \\
        \gamma(3t - 2) & \frac{2}{3} \leq t \leq 1
      \end{cases}
    \]
    Let
    \[
      f_0(t) =
      \begin{cases}
        \frac{4}{3}t & t \leq \frac{1}{2} \\
        \frac{1}{3} + \frac{2}{3} t & t \geq \frac{1}{2}
      \end{cases}
    \]
    and
    \[
      f_1(t) =
      \begin{cases}
        \frac{2}{3}t & t \leq \frac{1}{2} \\
        -\frac{1}{3} + \frac{4}{3}t & t \geq \frac{1}{2}
      \end{cases}
    \]
    Note that \(f_0 \simeq f\) as \emph{paths} via the straightline homotopy in \(I\). But
    \begin{align*}
      (\alpha \cdot \beta) \cdot \gamma &= \delta \compose f_0 \\
      \alpha \cdot (\beta \cdot \gamma) &= \delta \compose f_1
    \end{align*}
    so they are homotopic as path.
  \end{proof}
\end{proof}

\begin{eg}
  Let \(X = \R^n, x_0 = 0\). Consider a loop \(\gamma\) in \(\R^n\) based at \(0\). The straightline homotopy shows that \(\gamma \simeq c_0\) as path. Therefore \(\pi_1(\R^n, 0) \cong 1\).
\end{eg}

\begin{lemma}
  Let \(f: X \to Y\) be such that \(f(x_0) = y_0\). There is a well-defined homomorphism
  \begin{align*}
    f_*: \pi_1(X, x_0) &\to \pi_1(Y, y_0) \\
    [\gamma] &\mapsto [f \compose \gamma]
  \end{align*}
  Furthermore,
  \begin{enumerate}
  \item if \(f \simeq f' \rel \{x_0\}\) then \(f_* = f_*'\).
  \item if \(g: Y \to Z\) is another map then \(f_* \compose g_* = (f \compose g)_*\).
  \item \((\id_X)_* = \id_{\pi_1(X, x_0)}\).
  \end{enumerate}
\end{lemma}

\begin{proof}
  Easy.
\end{proof}

We'd like to eliminate the dependence of \(\pi_1(X, x_0)\) on \(x_0\), at least when \(X\) is path-connected. Suppose \(x_0, x_1 \in X\). What do \(\pi_1(X, x_0)\) and \(\pi_1(X, x_1)\) have to do with each other, where \(X\) is path-connected?

Fix \(\alpha\) a path from \(x_0\) to \(x_1\).

\begin{lemma}
  There is a well-defined group homomorphism
  \begin{align*}
    \alpha_\#: \pi_1(X, x_0) &\to \pi_1(X, x_1) \\
    [\gamma] &\mapsto [\overline \alpha \cdot \gamma \cdot \alpha]
  \end{align*}
  Furthermore
  \begin{enumerate}
  \item if \(\alpha \simeq \alpha'\) then \(\alpha_\# = \alpha_\#'\),
  \item \((c_{x_0})_\# = \id_{\pi_1(X, x_0)}\),
  \item if \(\beta\) is a path from \(x_1\) to \(x_2\), \(\beta_\# \compose \alpha_\# = (\alpha \cdot \beta)_\#\).
  \item if \(f: X \to Y\) then \((f \compose \alpha)_\# \compose f_* = f_* \compose \alpha_\#\).
  \end{enumerate}
\end{lemma}

Now it makes sense to talk about isomorphism type of the fundamental group of a path-connected space.

\begin{definition}[simply connected]\index{simply connected}
  If \(X\) is path-connected and \(\pi_1(X, x_0) \cong 1\) for some (i.e.\ any) \(x_0 \in X\) then we say \(X\) is \emph{simply connected}.
\end{definition}

Our last task is to understand what homotopies that don't fix basepoints do to the fundamental group.

\begin{lemma}
  Suppose \(f, g: X \to Y\) is such that \(f \simeq_F g\). Define \(\alpha(t) = F(x_0, t)\), a path from \(f(x_0)\) to \(g(x_0)\). Then the following diagram commutes:
\[
  \begin{tikzcd}
    & \pi_1(Y, f(x_0)) \ar[dd, "\alpha_\#"] \\
    \pi_1(X, x_0) \ar[ur, "f_*"] \ar[dr, "g_*"] \\
    & \pi_1(Y, g(x_0))
  \end{tikzcd}
\]
i.e.\ \(g_* = \alpha_\# \compose f_*\).
\end{lemma}

\begin{proof}
  Let \([\gamma] \in \pi_1(X, x_0)\). We need to show that
  \[
    [g \compose \gamma] = g_*[\gamma] = \alpha_\# \compose f_*[\gamma] = [\overline \alpha \cdot (f \compose \gamma) \cdot \alpha]
  \]
  which is saying
  \[
    g \compose \gamma \simeq \overline \alpha \cdot (f \compose \gamma) \cdot \alpha
  \]
  as paths. Consider
  \begin{align*}
    I \times I &\to Y \\
    (s, t) &\mapsto F(\gamma(s), t)
  \end{align*}
  Let \(H\) be the straightline homotopy in \(I \times I\) between the yellow path and the brown path. Then \(G \compose H\) is the homotopy we need.
\end{proof}

\begin{theorem}
  If \(f: X \to Y, g: Y \to X\) is a pair of homotopy equivalences and \(x_0 \in X\) then \(f_*: \pi_1(X, x_0) \to \pi_1(Y, f(x_0))\) is an isomorphism.
\end{theorem}

\begin{proof}
  Suffices to prove that \(f_*\) is bijective. Let \(g \compose f \simeq_F \id_X\) and \(\alpha\) be the path defined from \(F\) as above. Then
  \[
    g_* \compose f_* = (g \compose f)_* = \alpha_\# \compose \id_{\pi_1(X, x_0)} = \alpha_\#
  \]
  so \(f_*\) is injective. Similarly it is surjective.
\end{proof}































\printindex
\end{document}

% https://www.dpmms.cam.ac.uk/~hjrw2/teaching.html
